% BOOK-FILE-ID: f-5b7ced861c09
% section: 補助数列

\section{補助数列}\label{節：補助数列}
  \sectionref*{節：定数型}では,$b_n=g(n)a_n$という新しい数列を置いた.
  元の数列$a_n$のままでは,各項がどのような値になるのかすぐには見えない.
  しかし,$b_n$に目を移せば,すべての項が$b_1$に等しいことがすぐに分かった.
  そして最後に,$a_n=\frac{b_1}{g(n)}$として元の数列へ戻したのである.

  このように,元の数列を調べるために新しく置いた数列を\textbf{\bookterm{auxiliary}{補助数列}}という.
  補助数列を使う解法は,どれも次の3つの段階で進む.
  \begin{enumerate}
    \item 元の数列から補助数列を置く
    \item 補助数列について,既知の形を使って一般項を求める
    \item 置いたときの式を使って,元の数列へ戻す
  \end{enumerate}

  実は,補助数列はここで初めて現れたわけではない.
  \sectionref*{節：階差型}を思い出してほしい.
  そこでは,$a_{n+1}=a_n+f(n)$の右辺にある$f(n)$を,$f(1)$から$f(n-1)$まで足し合わせた.
  この$f(n)$は,隣り合う2項の差$a_{n+1}-a_n$に等しい.
  つまり,$b_n=a_{n+1}-a_n$と置いた数列,すなわち\bookterm{difference}{階差}数列を調べていたことになる.
  階差数列$b_n$の一般項は$f(n)$そのものであり,その和をとって$a_n=a_1+\sum_{k=1}^{n-1}b_k$として元の数列へ戻した.
  したがって,階差型で考えた階差数列も,補助数列の一つと見ることができる.
  先ほどの3つの段階に当てはめれば,$b_n=a_{n+1}-a_n$と置くことが1つ目,$b_n=f(n)$と読み取ることが2つ目,和をとって$a_n$へ戻すことが3つ目の段階にあたる.

  \sectionref{節：応用パターン}からは,様々な補助数列が登場する.
  例えば$b_n=a_n-\alpha$のように定数を引いたり,$b_n=\frac{1}{a_n}$のように逆数をとったりして,基本パターンの形を作るのだ.
  どの場合も,何を補助数列として置いたのか,そしてそれをどう元の数列へ戻すのかを意識しながら読み進めよう.

  \begin{bookpoint}
    \textbf{POINT}\par
    補助数列は,元の数列を調べるために新しく置いた数列である.
    補助数列で既知の形を作って解き,置いたときの式を使って元の数列へ戻す.
  \end{bookpoint}
